\begin{pastexam}[label={px:ust-115-ee-ds-44}]{臺聯大電機類 115 學年度 第 44 題}
（單選；2 分）

Given the following pseudo code of Quicksort algorithm, which of the following statement is the
\textbf{least accurate}?
\begin{minted}[linenos, numbersep=1.5mm, xleftmargin=2em, fontsize=\small]{text}
Quicksort(A,p,r)
    if p<r
        q=Partition(A,p,r)
        Quicksort(A,p,q-1)
        Quicksort(A,q+1,r)

Partition(A,p,r)
    x=A[r]
    i=p-1
    for j=p to r-1
        if A[j]<=x
            i=i+1
            exchange A[i] with A[j]
    exchange A[i+1] with A[r]
    return i+1
\end{minted}
\begin{enumerate}[label=(\Alph*)]
\item It is stable.
\item It is in-place.
\item Its worst time performance is $O(n^2)$.
\item The output of Quicksort() is a non-decreasing sorted sequence.
\item The variable \textbf{x} is usually called the pivot.
\end{enumerate}
\end{pastexam}

\begin{pastsoln}
答案是 (A)。

這是 Lomuto 分割的快速排序（資結書第~14.5~節，虛擬碼~14.8、14.9）。

\textbf{反例}：$A = \langle 2_a, 2_b, 1\rangle$（下標標出兩個 $2$ 原來的先後）。
\proc{Partition}$(A, 1, 3)$：$x = 1$；$j = 1, 2$ 時 $2 \le 1$ 都不成立，$i$ 停在 $0$；最後交換
$A[1]$ 和 $A[3]$，得 $\langle 1, 2_b, 2_a\rangle$，回傳 $q = 1$。 右半 \proc{Partition}$(A, 2, 3)$：
$x = 2_a$；$j = 2$ 時 $2_b \le 2_a$，$i = 2$，自己和自己交換；最後交換 $A[3]$ 和自己。 結果
$\langle 1, 2_b, 2_a\rangle$，兩個 $2$ 的先後顛倒了，不穩定。 原因是分割時的長距離交換（第一步把
$2_a$ 直接換到最後面）會跨過相等的元素。

逐項檢查：
\begin{itemize}
\item (A) 不穩定，見反例。 \textbf{最不準確}，選它。
\item (B) 只在陣列裡交換，除了遞迴堆疊不用額外陣列，是原地排序（遞迴堆疊平均 $O(\log n)$，最壞
$O(n)$）。 正確。
\item (C) 每次基準都是最大或最小值時（例如已經排好），分割成 $n - 1$ 和 $0$，
$T(n) = T(n-1) + \Theta(n) = \Theta(n^2)$。 正確。
\item (D) 第 11 行用 $\le$，分割後 $A[p..q-1] \le x \le A[q+1..r]$，遞迴排好兩邊後整段由小到大
（非遞減）。 正確。
\item (E) $x = A[r]$ 是拿來比較、分割的那個元素，就是基準（pivot）。 正確。
\end{itemize}
\end{pastsoln}
